\documentclass[11pt]{article}
\usepackage{amsmath,amssymb,amsthm,hyperref}
\newtheorem{theorem}{Theorem}
\newtheorem{corollary}{Corollary}
\title{A quadratic barrier for corrections in one column per row}
\author{Asymptotic research note; independently reviewed}
\date{30 September 2026}
\begin{document}\maketitle
The polynomial exceptional-die construction starts with distinct scalar
nodes and corrects at most one comparison column at each node. The
following argument shows that this particular strategy cannot reach a
subquadratic number of rows. It is not a general lower bound for dice.

\begin{theorem}
Let $m\ge3$ and let $t_1,\ldots,t_K\in[0,1]$ be pairwise distinct.
Suppose a $K\times m$ array has the form
\[
 p_j(q)=t_q+\delta_j(q),
\]
where at every row at most one $\delta_j(q)$ is nonzero, and suppose
\[
 \frac1K\sum_q\prod_{j\in S}p_j(q)=\frac1{|S|+1}
 \qquad(S\subseteq[m]).
\]
Then $K\ge c m^2$ for an absolute positive constant $c$.
More specifically, either the scalar nodes themselves give an
equal-weight interval quadrature of degree $m-1$, or at least
$m^2/2$ rows have a nonzero correction.
\end{theorem}
\begin{proof}
Put
\[
 b_{j,r}=\frac1K\sum_q\delta_j(q)t_q^r.
\]
At a row, all terms involving two corrections vanish. For a set $S$
of size $s$ the mixed-moment equation is therefore
\[
 \frac1K\sum_q t_q^s+\sum_{j\in S}b_{j,s-1}=\frac1{s+1}.
\]
For $1\le s\le m-1$, comparing two sets differing in one index shows
that $b_{j,s-1}$ is independent of $j$. Denote their common value by
$b_{s-1}$.

If all $b_0,\ldots,b_{m-2}$ vanish, the scalar nodes integrate every
monomial through degree $m-1$. The classical lower bound for equal-weight
uniform-interval quadrature gives $K\ge c m^2$.

Otherwise let $Q_j=\{q:\delta_j(q)\ne0\}$ and $k_j=|Q_j|$.
These sets are pairwise disjoint. For any two indices $i\ne j$ the
signed measure
\[
 \sum_{q\in Q_i}\delta_i(q)\delta_{t_q}
 -\sum_{q\in Q_j}\delta_j(q)\delta_{t_q}
\]
annihilates all polynomials of degree at most $m-2$.
Its support consists of $k_i+k_j$ distinct points and all its
coefficients are nonzero. If $k_i+k_j\le m-1$, the Vandermonde matrix
would force every coefficient to vanish, which would also force all
the common $b_r$ to vanish. Consequently $k_i+k_j\ge m$ for every
pair. Summing over pairs gives
\[
 (m-1)\sum_j k_j\ge m\binom m2,
 \qquad K\ge\sum_jk_j\ge\frac{m^2}{2}.
\]
\end{proof}

The common correction functional is forced only through degree $m-2$;
the full $m$-column equation supplies one additional aggregate condition.
The proof deliberately does not assume equality of the degree-$m-1$
correction moments. Pairwise distinctness of the baseline nodes is an
explicit hypothesis. The result applies to the separated-node construction
in \path{size_bounds/polynomial_exceptional_die.tex}, but not to an
arbitrary monotone comparison tensor.

For the classical scalar quadrature bound, see the uniform-weight case
in Gilboa--Peled, \emph{Chebyshev-type quadratures for doubling weights},
\url{https://arxiv.org/abs/1507.01505}.
\begin{corollary}[Bounded overlap]
Let $t_1,\ldots,t_K\in[0,1]$ be distinct, and suppose
$p_j(q)=t_q+\delta_j(q)$, $1\le j\le m$, has uniform-interval mixed moments
\[
 K^{-1}\sum_q\prod_{j\in S}p_j(q)=\frac1{|S|+1}.
\]
Suppose each row has nonzero corrections in at most $d\ge1$ columns.
There is an absolute $c>0$ such that
\[
 K\ge c\min\left\{m^2,\frac{m^{4/3}}{d^{4/3}}\right\}.
\]
In particular any fixed bound on the number of corrected columns per
row requires $K=\Omega(m^{4/3})$. For $d=1$ the stronger bound
$K=\Omega(m^2)$ holds.
\end{corollary}
\begin{proof}
Join two columns by an edge if some row corrects both of them. This
conflict graph has at most $K\binom d2$ edges. It has an independent
set of size at least
\[
 R\ge\frac{m^2}{m+K d(d-1)}.
\tag{1}
\]
For completeness, order the vertices uniformly at random and retain
those appearing before every neighbor. The retained set is independent,
its expected size is $\sum_j(1+\deg j)^{-1}$, and Cauchy--Schwarz gives
(1).

On the retained columns each row corrects at most one column, and all
mixed-moment identities persist. The one-column theorem therefore gives
$K\ge c_0R^2$ for an absolute $c_0>0$. This estimate can be made valid
also for $R=1,2$ by reducing $c_0$, since $K\ge1$.
If $Kd(d-1)\le m$, (1) gives $K\ge c_0m^2/4$.
Otherwise (1) and the same estimate give
\[
 K^3[d(d-1)]^2\ge c_0m^4/4.
\]
The claimed bound follows. When $d=1$ use the one-column theorem directly.
\end{proof}


\end{document}
